%% ============================================================
%% SPRINT 6 — EXPLICIT TABLE OF CONTENTS PREVIEW
%% Physics-Informed Game-Theoretic Defense of Pipeline Infrastructure
%%
%% Purpose : Serves as the document map for Sprint 6 writing.
%%           The actual compiled TOC is auto-generated by \tableofcontents
%%           in thesis.tex; this file provides a human-readable structural
%%           overview for review, supervision meetings, and chapter planning.
%%
%% Status  : All sections confirmed present in chapter .tex files.
%%           Word-count targets are estimates for a ~25 000-word MSc thesis.
%%
%% Author  : Babak Pirzadi  |  STRATEGOS MSc
%% Updated : Sprint 6 kickoff
%% ============================================================

\documentclass[12pt, a4paper]{article}
\usepackage[top=2cm, bottom=2cm, left=2.5cm, right=2.5cm]{geometry}
\usepackage[T1]{fontenc}
\usepackage[utf8]{inputenc}
\usepackage{mathptmx}
\usepackage{microtype}
\usepackage[dvipsnames]{xcolor}
\usepackage{booktabs}
\usepackage{array}
\usepackage{longtable}
\usepackage{setspace}
\usepackage{titlesec}
\usepackage{hyperref}

\definecolor{navyblue}{RGB}{26,58,92}
\definecolor{safegreen}{RGB}{30,132,73}
\definecolor{strategyred}{RGB}{192,57,43}
\definecolor{lightgray}{RGB}{245,245,245}

\hypersetup{colorlinks=true, linkcolor=navyblue, urlcolor=strategyred}

\titleformat{\section}{\bfseries\large\color{navyblue}}{\thesection}{1em}{}
\titleformat{\subsection}{\bfseries\normalsize\color{navyblue}}{\thesubsection}{1em}{}

\begin{document}

\begin{center}
  {\LARGE\bfseries\color{navyblue}
    Physics-Informed Game-Theoretic Defense\\[0.3em]
    of Pipeline Infrastructure}\\[0.6em]
  {\large\color{navyblue}
    Integrating Weld Joint Vulnerability into\\
    Bayesian Stackelberg Attacker--Defender Models}\\[0.8em]
  \rule{0.6\textwidth}{0.6pt}\\[0.4em]
  {\normalsize Babak Pirzadi \quad|\quad STRATEGOS MSc}\\[0.2em]
  {\small\textit{Sprint 6 — Document Map \& Table of Contents Preview}}
\end{center}

\vspace{1em}

%% ============================================================
%%  STRUCTURAL OVERVIEW
%% ============================================================
\section*{\color{navyblue}Thesis Structure at a Glance}

\begin{longtable}{@{}p{0.08\textwidth} p{0.52\textwidth} p{0.16\textwidth} p{0.12\textwidth}@{}}
\toprule
\textbf{Item} & \textbf{Title} & \textbf{Target words} & \textbf{Status} \\
\midrule
\endhead
\bottomrule
\endfoot

\multicolumn{4}{@{}l}{\textit{Front Matter}} \\[2pt]
--- & Declaration of Originality & --- & \textcolor{safegreen}{Ready} \\
--- & Abstract & 300 & \textcolor{safegreen}{Ready} \\
--- & Acknowledgements & 150 & \textcolor{safegreen}{Ready} \\
--- & List of Abbreviations & --- & \textcolor{safegreen}{Ready} \\
--- & List of Figures (Figs.\ 1--20) & --- & auto \\
--- & List of Tables & --- & auto \\[6pt]

\multicolumn{4}{@{}l}{\textit{Main Chapters}} \\[2pt]
Ch.\ 1 & Introduction & 2\,000 & \textcolor{safegreen}{Draft} \\
Ch.\ 2 & Literature Review & 3\,500 & \textcolor{safegreen}{Draft} \\
Ch.\ 3 & Physics-Informed Vulnerability Engine (Zone C) & 5\,500 & \textcolor{safegreen}{Draft} \\
Ch.\ 4 & Adversarial Robustness of NDE Classifiers (Zone A) & 4\,500 & \textcolor{safegreen}{Draft} \\
Ch.\ 5 & Results and Integrated Analysis & 4\,000 & \textcolor{safegreen}{Draft} \\
Ch.\ 6 & Conclusion and Future Work & 2\,000 & \textcolor{safegreen}{Draft} \\[6pt]

\multicolumn{4}{@{}l}{\textit{End Matter}} \\[2pt]
--- & Bibliography & --- & \textcolor{safegreen}{Ready} \\
App.\ A & LP Derivation & 800 & pending \\
App.\ B & FAD Reference Tables & 500 & pending \\
App.\ C & Code Listings (key functions) & 600 & pending \\
\bottomrule
\end{longtable}

\vspace{1.2em}

%% ============================================================
%%  DETAILED TABLE OF CONTENTS
%% ============================================================
\section*{\color{navyblue}Detailed Table of Contents}
\onehalfspacing

%% ---- FRONT MATTER -------------------------------------------
\noindent{\bfseries\color{navyblue}Front Matter}\vspace{4pt}

\noindent
\begin{tabular}{@{}p{0.04\textwidth}p{0.78\textwidth}p{0.1\textwidth}@{}}
  & Declaration of Originality & \textit{i} \\
  & Abstract & \textit{ii} \\
  & Acknowledgements & \textit{iii} \\
  & List of Abbreviations & \textit{iv} \\
  & List of Figures & \textit{vi} \\
  & List of Tables & \textit{viii} \\
\end{tabular}

\vspace{0.8em}

%% ---- CHAPTER 1 -----------------------------------------------
\noindent{\bfseries\large\color{navyblue}1\quad Introduction}\hfill\textit{p.\ 1}

\vspace{4pt}
\noindent
\begin{tabular}{@{}p{0.08\textwidth}p{0.78\textwidth}p{0.06\textwidth}@{}}
  1.1 & Motivation and Strategic Context &  \\
  & \small\textit{TEN-E network; Colonial Pipeline \& Nord Stream as case primes;} & \\
  & \small\textit{three-domain gap: structural integrity \(\cap\) network resilience \(\cap\) adversarial intent.} & \\[4pt]
  1.2 & Problem Statement & \\
  & \small\textit{Decoupling of physical vulnerability from game-theoretic payoffs;} & \\
  & \small\textit{formalisation of the physics-to-utility mapping objective.} & \\[4pt]
  1.3 & Thesis Contributions & \\
  & \small\textit{Five enumerated contributions: FAD–SSE coupling; PHMSA calibration;} & \\
  & \small\textit{Bayesian Stackelberg LP; adversarial NDE threat model; Dash dashboard.} & \\[4pt]
  1.4 & Thesis Scope and Limitations & \\
  & \small\textit{Synthetic Gulf Coast network; U.S.\ PHMSA data; white-box attack assumption.} & \\[4pt]
  1.5 & Regulatory Framework & \\
  & \small\textit{ASME B31.8; BS 7910:2019; API 1104; API 579; PHMSA Part 192.} & \\[4pt]
  1.6 & Thesis Structure & \\
  & \small\textit{Chapter-by-chapter roadmap.} & \\
\end{tabular}

\vspace{0.8em}

%% ---- CHAPTER 2 -----------------------------------------------
\noindent{\bfseries\large\color{navyblue}2\quad Literature Review}\hfill\textit{p.\ 15}

\vspace{4pt}
\noindent
\begin{tabular}{@{}p{0.08\textwidth}p{0.78\textwidth}p{0.06\textwidth}@{}}
  2.1 & Overview & \\
  & \small\textit{Taxonomy of related work across three pillars.} & \\[4pt]
  2.2 & Pipeline Structural Integrity and Failure Assessment & \\[2pt]
  2.2.1 & \quad Fracture Mechanics and Fitness-for-Service & \\
  & \quad\small\textit{LEFM; \(K_r/L_r\) failure assessment diagram; BS 7910 vs.\ API 579 approaches.} & \\[2pt]
  2.2.2 & \quad Fatigue Design of Welded Joints & \\
  & \quad\small\textit{IIW S-N classification; Miner's rule; weld class hierarchy (D–W).} & \\[2pt]
  2.2.3 & \quad Monte Carlo Methods in Reliability Assessment & \\
  & \quad\small\textit{Lognormal SCF variability; 10 000-simulation convergence; literature precedent.} & \\[4pt]
  2.3 & Security Games for Critical Infrastructure & \\[2pt]
  2.3.1 & \quad Stackelberg Security Games: Foundations & \\
  & \quad\small\textit{Leader-follower; DOBSS; SSE uniqueness; Conitzer \& Sandholm 2006.} & \\[2pt]
  2.3.2 & \quad Deployed Systems & \\
  & \quad\small\textit{ARMOR (LAX); PROTECT (ferries); IRIS (air marshals).} & \\[2pt]
  2.3.3 & \quad Critical Infrastructure and Network Attacks & \\
  & \quad\small\textit{Graph-theoretic attack models; betweenness-centrality targeting; max-flow disruption.} & \\[4pt]
  2.4 & Adversarial Machine Learning and NDE Security & \\[2pt]
  2.4.1 & \quad Adversarial Attacks on Neural Networks & \\
  & \quad\small\textit{Szegedy 2014; FGSM (Goodfellow 2015); BIM (Kurakin 2017); PGD (Madry 2018).} & \\[2pt]
  2.4.2 & \quad Machine Learning for NDE & \\
  & \quad\small\textit{CNN/MLP weld inspection; ultrasonic feature classification; sparse literature on adversarial NDE.} & \\[4pt]
  2.5 & The Unifying Gap & \\
  & \small\textit{No prior work couples BS 7910 \(P_f\) directly to SSE utilities; no adversarial NDE threat model} & \\
  & \small\textit{linked to game-theoretic consequences — the gap this thesis fills.} & \\
\end{tabular}

\vspace{0.8em}

%% ---- CHAPTER 3 -----------------------------------------------
\noindent{\bfseries\large\color{navyblue}3\quad Physics-Informed Vulnerability Engine (Zone C)}\hfill\textit{p.\ 30}

\vspace{4pt}
\noindent
\begin{tabular}{@{}p{0.08\textwidth}p{0.78\textwidth}p{0.06\textwidth}@{}}
  3.1 & Chapter Overview & \\
  & \small\textit{Zone C architecture; four-module pipeline (calibration \(\to\) FAD \(\to\) MC \(\to\) network \(\to\) game).} & \\[4pt]
  3.2 & PHMSA Empirical Calibration & \\[2pt]
  3.2.1 & \quad Dataset & \\
  & \quad\small\textit{1\,996 incident records, 2010–2024; 216 material failure events; seam-type stratification.} & \\[2pt]
  3.2.2 & \quad Calibrated Distributions & \\
  & \quad\small\textit{SCF hierarchy: seamless 1.30, ERW-HF 1.50, ERW-LF 1.70, DSAW 1.45, fillet 2.10;} & \\
  & \quad\small\textit{SMYS lognormal \(\mu=448\) MPa, \(\sigma=0.04\); fracture toughness \(\mu=55\) MPa\(\sqrt{\text{m}}\), \(\sigma=0.12\).} & \\[4pt]
  3.3 & BS 7910 Failure Assessment Diagram Engine & \\[2pt]
  3.3.1 & \quad Option 1 FAD Curve & \\
  & \quad\small\textit{\(f(L_r) = (1 - 0.14L_r^2)[0.3 + 0.7\exp(-0.65L_r^6)]\); implementation in \texttt{fad\_engine.py}.} & \\[2pt]
  3.3.2 & \quad Toughness and Load Ratios & \\
  & \quad\small\textit{\(K_r = K_I/K_{mat}\); \(L_r = P/P_Y\); \(L_{r,\max} = 1.15\) (X60), \(1.12\) (X65); assess\_flaw() output.} & \\[4pt]
  3.4 & Monte Carlo Failure Probability Engine & \\[2pt]
  3.4.1 & \quad Simulation Design & \\
  & \quad\small\textit{10\,000 simulations/segment; seed = 42; lognormal SCF, SMYS, toughness sampling;} & \\
  & \quad\small\textit{\(P_f\) = fraction of realisations in which assessment point falls outside FAD envelope.} & \\[4pt]
  3.5 & Gulf Coast Network Model & \\[2pt]
  3.5.1 & \quad Topology & \\
  & \quad\small\textit{20 nodes (1 source, 2 compressors, 8 delivery, 3 storage, 6 junctions);} & \\
  & \quad\small\textit{22 directed segments; betweenness centrality range [0.00, 0.26].} & \\[2pt]
  3.5.2 & \quad \(P_f\) Attachment & \\
  & \quad\small\textit{\(P_f\) range [0.29, 0.93]; top-risk segment SEG\_N012\_N015 (\(P_f = 0.93\), max centrality).} & \\[4pt]
  3.6 & Bayesian Stackelberg Security Game & \\[2pt]
  3.6.1 & \quad Game Formulation & \\
  & \quad\small\textit{Defender \(\bm{c} \in [0,1]^{22}\); attacker best-response; Bayesian payoff over 3 type profiles.} & \\[2pt]
  3.6.2 & \quad LP Formulation and Strong Stackelberg Equilibrium & \\
  & \quad\small\textit{Milind-type LP enumeration; \(\sum c_i \leq B\); \(\delta = 0.25\) interception bonus.} & \\[2pt]
  3.6.3 & \quad Attacker Type Profiles & \\
  & \quad\small\textit{STRATEGIC (prior 0.50): rational \(P_f\)-maximiser;} & \\
  & \quad\small\textit{OPPORTUNISTIC (0.30): targets by network degree;} & \\
  & \quad\small\textit{STATE\_ACTOR (0.20): exploits both \(P_f\) and centrality.} & \\[2pt]
  3.6.4 & \quad Game Results & \\
  & \quad\small\textit{SSE defender utility: \(-0.477\); coverage effectiveness vs.\ zero-budget: \(+52.3\)\%;} & \\
  & \quad\small\textit{risk reduction vs.\ uniform: \(\mathbf{17.0\%}\); diminishing returns beyond \(\approx 60\%\) budget.} & \\
\end{tabular}

\vspace{0.8em}

%% ---- CHAPTER 4 -----------------------------------------------
\noindent{\bfseries\large\color{navyblue}4\quad Adversarial Robustness of NDE Classifiers (Zone A)}\hfill\textit{p.\ 58}

\vspace{4pt}
\noindent
\begin{tabular}{@{}p{0.08\textwidth}p{0.78\textwidth}p{0.06\textwidth}@{}}
  4.1 & Chapter Overview & \\
  & \small\textit{Threat narrative: state-actor suppresses NDE signal \(\Rightarrow\) hides defect \(\Rightarrow\) exploits physics gap.} & \\[4pt]
  4.2 & The Adversarial NDE Threat & \\
  & \small\textit{Cross-layer attack vector: Zone A (inspection) \(\to\) Zone C (physics) \(\to\) Zone B (cyber) coupling.} & \\[4pt]
  4.3 & Synthetic NDE Feature Dataset & \\[2pt]
  4.3.1 & \quad Feature Engineering & \\
  & \quad\small\textit{32 features: amplitude \(\times 8\), frequency \(\times 8\), geometry \(\times 8\), texture \(\times 8\);} & \\
  & \quad\small\textit{4 classes: CLEAN (0), POROSITY (1), CRACK (2), LACK\_OF\_FUSION (3); noise\_level = 0.18.} & \\[2pt]
  4.3.2 & \quad Dataset Generation & \\
  & \quad\small\textit{Stratified synthetic generation; class-conditional Gaussian clusters; PHMSA-scaled severity.} & \\[4pt]
  4.4 & WeldDefectMLP Architecture & \\[2pt]
  4.4.1 & \quad Architecture & \\
  & \quad\small\textit{\(32 \to 128 \to 64 \to 4\); Xavier initialisation; ReLU activations; Dropout(0.30); Softmax output.} & \\[2pt]
  4.4.2 & \quad Training Procedure & \\
  & \quad\small\textit{80 epochs; cosine LR annealing (0.05 \(\to\) 10\textsuperscript{-4}); SGD + momentum = 0.9;} & \\
  & \quad\small\textit{clean accuracy 99.7\%; numerical gradient check \(<5\%\) relative error.} & \\[4pt]
  4.5 & Adversarial Attacks & \\[2pt]
  4.5.1 & \quad Fast Gradient Sign Method (FGSM) & \\
  & \quad\small\textit{Single-step \(L_\infty\): \(\delta = \varepsilon \cdot \text{sign}(\nabla_x \mathcal{L})\); Goodfellow et al.\ 2015.} & \\[2pt]
  4.5.2 & \quad Basic Iterative Method (BIM) & \\
  & \quad\small\textit{Iterative FGSM without random restart; 10 steps, step size \(\alpha = \varepsilon/10\); Kurakin et al.\ 2017.} & \\[2pt]
  4.5.3 & \quad Projected Gradient Descent (PGD) & \\
  & \quad\small\textit{Madry et al.\ 2018; random \(\delta\)-ball initialisation; projected onto \(L_\infty(\varepsilon)\) ball at each step.} & \\[2pt]
  4.5.4 & \quad Physics-Informed Epsilon Scaling & \\
  & \quad\small\textit{\(\varepsilon_{\mathrm{eff}} = 0.30 \times \text{SCF}/1.5\); segment-level adversarial threat score in dashboard.} & \\[4pt]
  4.6 & Results & \\[2pt]
  4.6.1 & \quad Classification Performance & \\
  & \quad\small\textit{Clean accuracy 99.7\%; confusion matrices (Figs.\ 17a--17d); per-class F1 reported.} & \\[2pt]
  4.6.2 & \quad Attack Success Rates & \\
  & \quad\small\textit{At \(\varepsilon = 0.30\): FGSM ASR 18.1\% (acc.\ \(\to\) 81.7\%); BIM ASR 20.1\% (acc.\ \(\to\) 79.7\%);} & \\
  & \quad\small\textit{PGD ASR 9.4\% (acc.\ \(\to\) 90.3\%); epsilon-sweep monotonicity validated (Fig.\ 16).} & \\[2pt]
  4.6.3 & \quad The BIM--PGD Gap: Interpretation & \\
  & \quad\small\textit{BIM > PGD: local gradient curvature dominates; random initialisation loses in feature domain;} & \\
  & \quad\small\textit{implication: adversarial training should target BIM-generated examples for NDE hardening.} & \\
\end{tabular}

\vspace{0.8em}

%% ---- CHAPTER 5 -----------------------------------------------
\noindent{\bfseries\large\color{navyblue}5\quad Results and Integrated Analysis}\hfill\textit{p.\ 80}

\vspace{4pt}
\noindent
\begin{tabular}{@{}p{0.08\textwidth}p{0.78\textwidth}p{0.06\textwidth}@{}}
  5.1 & Chapter Overview & \\
  & \small\textit{Integrated reading across Zone C + Zone A layers; Dashboard as unified evidence platform.} & \\[4pt]
  5.2 & Network Vulnerability Landscape & \\
  & \small\textit{Per-segment \(P_f\) heat-map; seam-type stratification; betweenness--\(P_f\) correlation analysis (Fig.\ 11).} & \\[4pt]
  5.3 & Benefit of Physics-Informed Stackelberg Allocation & \\[2pt]
  5.3.1 & \quad Scenario Definitions & \\
  & \quad\small\textit{Scenario A: zero-budget (baseline); Scenario B: uniform coverage; Scenario C: SSE.} & \\[2pt]
  5.3.2 & \quad Total Expected Network Loss & \\
  & \quad\small\textit{Bayesian utility across 3 attacker types; SSE achieves \(\mathbf{-17.0\%}\) vs.\ uniform baseline;} & \\
  & \quad\small\textit{\(+52.3\%\) effectiveness vs.\ zero-budget; budget sensitivity (Fig.\ 14).} & \\[4pt]
  5.4 & Adversarial NDE and the Game-Theoretic Attacker & \\
  & \small\textit{Cross-layer threat integration: FGSM/BIM/PGD ASR calibrated per segment via} & \\
  & \small\textit{\(\varepsilon_{\mathrm{eff}}\); STATE\_ACTOR scenario analysis; CRACK \(\to\) CLEAN misclassification consequence.} & \\[4pt]
  5.5 & Interactive Dashboard & \\
  & \small\textit{Tab 1: Network Intelligence — click-to-inspect FAD + adversarial gauge (Fig.\ 18);}& \\
  & \small\textit{Tab 2: Stackelberg Defender — SSE/baseline heatmap + budget slider (Fig.\ 19);}& \\
  & \small\textit{Tab 3: Scenario Comparison — residual risk bars + robustness curves (Fig.\ 20).} & \\[4pt]
  5.6 & Discussion & \\
  & \small\textit{Physics coupling as information asymmetry reducer; threshold budget effects;} & \\
  & \small\textit{adversarial NDE as intelligence preparation of the battlefield.} & \\
\end{tabular}

\vspace{0.8em}

%% ---- CHAPTER 6 -----------------------------------------------
\noindent{\bfseries\large\color{navyblue}6\quad Conclusion and Future Work}\hfill\textit{p.\ 100}

\vspace{4pt}
\noindent
\begin{tabular}{@{}p{0.08\textwidth}p{0.78\textwidth}p{0.06\textwidth}@{}}
  6.1 & Summary of Contributions & \\
  & \small\textit{Five contributions restated against the stated objectives of Ch.\ 1.3;} & \\
  & \small\textit{310-test validated codebase; 17.0\% risk reduction headline finding.} & \\[4pt]
  6.2 & Implications for Mediterranean Pipeline Security & \\
  & \small\textit{TAP corridor (Greece--Albania--Italy); SNAM Rete Gas distribution nodes;} & \\
  & \small\textit{ENTSOG 2023 data as calibration source; EU CER Directive 2022/2557 framing.} & \\[4pt]
  6.3 & Limitations & \\
  & \small\textit{Synthetic network (no real SCADA topology); white-box attack assumption;} & \\
  & \small\textit{single ILI inspection modality; Stackelberg rationality assumption.} & \\[4pt]
  6.4 & Directions for Future Research & \\
  & \small\textit{Real TAP/SNAM topology integration; grey-box \& transfer attacks on NDE;} & \\
  & \small\textit{multi-period Stackelberg; ILI sensor fusion; EU-CERT information-sharing protocol.} & \\[4pt]
  6.5 & Closing Remarks & \\
  & \small\textit{Convergence of structural integrity and security science as a strategic discipline.} & \\
\end{tabular}

\vspace{0.8em}

%% ---- END MATTER -----------------------------------------------
\noindent{\bfseries\large\color{navyblue}End Matter}

\vspace{4pt}
\noindent
\begin{tabular}{@{}p{0.08\textwidth}p{0.78\textwidth}p{0.06\textwidth}@{}}
  --- & Bibliography & \textit{p.\ 112} \\[4pt]
  App.\ A & LP Derivation: Milind-Type Enumeration for Bayesian SSE & \textit{p.\ 125} \\
  & \small\textit{Full LP variable set; dual feasibility proof sketch; Python scipy.optimize.linprog call.} & \\[4pt]
  App.\ B & FAD Reference Tables: Lr$_{\max}$ and Material Constants & \textit{p.\ 130} \\
  & \small\textit{Grade X52/X60/X65/X70 constants; SCF lookup table by seam type.} & \\[4pt]
  App.\ C & Code Listings: Core Algorithms & \textit{p.\ 133} \\
  & \small\textit{\texttt{assess\_flaw()}, \texttt{solve\_stackelberg\_game()}, \texttt{WeldDefectMLP.forward()},} & \\
  & \small\textit{\texttt{fgsm\_attack()}, \texttt{pgd\_attack()}.} & \\
\end{tabular}

\vspace{1.5em}

%% ============================================================
%%  FIGURE MAPPING TO CHAPTERS
%% ============================================================
\section*{\color{navyblue}Figure Assignment by Chapter}

\begin{longtable}{@{}p{0.08\textwidth} p{0.12\textwidth} p{0.65\textwidth}@{}}
\toprule
\textbf{Figure} & \textbf{Chapter} & \textbf{Caption theme (strategic framing)} \\
\midrule
\endhead
\bottomrule
\endfoot
Fig.\ 1  & Ch.\ 3 & BS 7910 FAD: nominal flaws within envelope; critical flaws outside for high-SCF seams \\
Fig.\ 2  & Ch.\ 3 & FAD comparison across X52--X70 grades: Lr$_{\max}$ sensitivity \\
Fig.\ 3  & Ch.\ 3 & IIW S-N curves: 2-order-of-magnitude fatigue life spread by weld class \\
Fig.\ 4  & Ch.\ 3 & MC \(P_f\) distribution: bimodal failure clustering by seam type \\
Fig.\ 5  & Ch.\ 3 & \(P_f\)--SCF sensitivity: linear amplification in the critical regime \\
Fig.\ 6  & Ch.\ 3 & PHMSA cause distribution: material failure as primary sabotage vector \\
Fig.\ 7  & Ch.\ 3 & PHMSA temporal trend 2010--2024: declining rate but stable material share \\
Fig.\ 8  & Ch.\ 3 & Calibrated parameter distributions: SCF, SMYS, toughness posteriors \\
Fig.\ 9  & Ch.\ 3 & SCF hierarchy: fillet joints \(\approx 1.6\times\) seamless pipe \\
Fig.\ 10 & Ch.\ 3 & Calibrated \(P_f\) by weld type: reproduces PHMSA failure-rate ordering \\
Fig.\ 11 & Ch.\ 3 & Network vulnerability map: betweenness vs.\ \(P_f\) for 22 segments \\
Fig.\ 12 & Ch.\ 3 & Network property distributions: diameter, wall thickness, operating pressure \\
Fig.\ 13 & Ch.\ 3 & SSE coverage probability map: disproportionate allocation to high-\(P_f\), high-centrality nodes \\
Fig.\ 14 & Ch.\ 3 & Budget sensitivity: diminishing returns beyond \(\approx 60\%\) allocation \\
Fig.\ 15 & Ch.\ 4 & Feature perturbation: clean vs.\ adversarial vs.\ \(\delta\) for CRACK class \\
Fig.\ 16 & Ch.\ 4 & Robustness curves: FGSM/BIM/PGD accuracy + ASR vs.\ \(\varepsilon\) \\
Fig.\ 17 & Ch.\ 4 & Confusion matrices: clean vs.\ post-attack misclassification patterns \\
Fig.\ 18 & Ch.\ 5 & Dashboard Tab 1 — Network Intelligence (click-to-inspect FAD + FGSM gauge) \\
Fig.\ 19 & Ch.\ 5 & Dashboard Tab 2 — Stackelberg coverage heatmap + budget slider \\
Fig.\ 20 & Ch.\ 5 & Dashboard Tab 3 — Scenario comparison + adversarial robustness curves \\
\end{longtable}

\vspace{1em}

%% ============================================================
%%  KEY NUMERICAL ANCHOR TABLE (for cross-chapter consistency)
%% ============================================================
\section*{\color{navyblue}Key Numerical Anchors — Single Source of Truth}
\small
\noindent All values sourced from \texttt{TECHNICAL\_INVENTORY.md}. Every table in the manuscript must agree with these numbers exactly.

\vspace{6pt}
\begin{longtable}{@{}p{0.36\textwidth} p{0.32\textwidth} p{0.18\textwidth}@{}}
\toprule
\textbf{Quantity} & \textbf{Value} & \textbf{Chapter} \\
\midrule
\endhead
\bottomrule
\endfoot
PHMSA records (2010--2024) & 1\,996 records, 216 material failures & Ch.\ 3 \\
Network topology & 20 nodes, 22 segments & Ch.\ 3 \\
\(P_f\) range (seed = 42, 5\,000 MC sims) & [0.29, 0.93] & Ch.\ 3 \\
\(P_f\): seamless / ERW-HF / ERW-LF / DSAW / fillet & 0.53 / 0.66 / 0.74 / 0.62 / 0.88 & Ch.\ 3 \\
FAD \(L_{r,\max}\) (X60 / X65) & 1.15 / 1.12 & Ch.\ 3 \\
Monte Carlo simulations / segment & 10\,000 (calibration); 5\,000 (network) & Ch.\ 3 \\
Attacker types \& priors & STRATEGIC 0.50; OPP 0.30; STATE 0.20 & Ch.\ 3 \\
Defender budget & 30\% (6.6 coverage units / 22) & Ch.\ 3 \\
SSE defender utility & \(-0.477\) & Ch.\ 3 \\
Coverage effectiveness vs.\ zero-budget & \(+52.3\%\) & Ch.\ 3, 5 \\
\textbf{Risk reduction vs.\ uniform (headline)} & \textbf{17.0\%} & Ch.\ 3, 5 \\
WeldDefectMLP architecture & \(32 \to 128 \to 64 \to 4\) & Ch.\ 4 \\
Clean accuracy & 99.7\% & Ch.\ 4 \\
FGSM ASR (\(\varepsilon = 0.30\)) & 18.1\% (accuracy: 81.7\%) & Ch.\ 4 \\
BIM ASR (\(\varepsilon = 0.30\)) & 20.1\% (accuracy: 79.7\%) & Ch.\ 4 \\
PGD ASR (\(\varepsilon = 0.30\)) & 9.4\% (accuracy: 90.3\%) & Ch.\ 4 \\
Numerical gradient check & \(<5\%\) relative error & Ch.\ 4 \\
\(\varepsilon_{\mathrm{eff}}\) scaling formula & \(0.30 \times \text{SCF} / 1.5\) & Ch.\ 4, 5 \\
Test suite & 310 tests passing, 0 failures & All \\
\end{longtable}

\end{document}
